\chapter*{孟子題辭}
\markboth{孟子題辭}{孟子題辭}
\addcontentsline{toc}{chapter}{孟子題辭}

孟子題辭者，所以題號孟子之書，本末指義，文辭之表也。孟，姓也；子者，男子之通稱也。此書，孟子之所作也，故總謂之《孟子》。其篇目則各自有名。

孟子，鄒人也，名軻，字則未聞也。鄒本春秋邾子之國，至孟子時改曰鄒矣。國近魯，後爲魯所幷。又言邾爲楚所幷，非魯也。今鄒縣是也。或曰：孟子，魯公族孟孫之後，故孟子仕於齊，喪母而歸葬於魯也。

三桓子孫旣以衰微，分適他國。孟子生有淑質，夙喪其父，幼被慈母三遷之教，長師孔子之孫子思，治儒術之道，通五經，尤長於《詩》《書》。

周衰之末，戰國縱橫，用兵爭彊以相侵奪。當世取士，務先權謀，以爲上賢。先王大道，陵遲墮廢，異端竝起。若楊朱、墨翟放蕩之言，以干時惑衆者非一。

孟子閔悼堯、舜、湯、文、周、孔之業將遂湮微，正塗壅底，仁義荒怠，佞僞馳騁，紅紫亂朱。於是則慕仲尼周流憂世，遂以儒道遊於諸侯，思濟斯民。然由不肻枉尺直尋，時君咸謂之迂闊於事，終莫能聽納其說。

孟子亦自知遭蒼姬之訖録，值炎劉之未奮。進不得佐興唐虞雍熙之和，退不能信三代之餘風，恥沒世而無聞焉。是故垂憲言以詒後人。仲尼有云：「我欲託之空言，不如載之行事之深切著明也。」

於是退而論集所與髙第弟子公孫丑、萬章之徒難疑答問，又自撰其法度之言，著書七篇，二百六十一章，三萬四千六百八十五字。包羅天地，揆敘萬類，仁義道德，性命禍福，粲然靡所不載。

帝王公侯遵之，則可以致隆平，頌淸廟；卿大夫士蹈之，則可以尊君父，立忠信；守志厲操者儀之，則可以崇髙節，抗浮雲。有風人之託物，二雅之正言，可謂直而不倨，曲而不屈，命世亞聖之大才者也。

孔子自衞反魯，然後樂正，《雅》《頌》各得其所。乃刪《詩》、定《書》、繫《周易》、作《春秋》。孟子退自齊梁，述堯舜之道而著作焉。此大賢擬聖而作者也。

七十子之疇，會集夫子所言，以爲《論語》。《論語》者，五經之錧鎋，六藝之喉衿也。《孟子》之書則而象之。

衞靈公問陳於孔子，孔子答以俎豆；梁惠王問利國，孟子對以仁義。宋桓魋欲害孔子，孔子稱「天生德於予」；魯臧倉毀鬲孟子，孟子曰「臧氏之子，焉能使予不遇哉？」旨意合同，若此者衆。

又有外書四篇：《性善》、《辯文》、《說孝經》爲正。其文不能弘深，不與內篇相似，似非孟子本眞，後世依放而託之者也。

孟子旣沒之後，大道遂絀。逮至亡秦，焚滅經術，坑戮儒生，孟子徒黨盡矣。其書號爲諸子，故篇籍得不泯絕。

漢興，除秦虐禁，開延道德。孝文皇帝欲廣遊學之路，《論語》、《孝經》、《孟子》、《爾雅》皆置博士。後罷傳記博士，獨立五經而已。

訖今諸經通義，得引《孟子》以明事，謂之博文。孟子長於譬喻，辭不迫切，而意已獨至。其言曰：「說《詩》者不以文害辭，不以辭害志，以意逆志，爲得之矣。」斯言殆欲使後人深求其意，以解其文，不但施於說《詩》也。

今諸解者徃徃摭取而說之，其說又多乖異不同。孟子以來五百餘載，傳之者亦已衆多。

余生西京，世尋丕祚，有自來矣。少蒙義方，訓渉典文。知命之際，嬰戚于天，遘屯離蹇，詭姓遁身，經營八紘之內，十有餘年。心勦形瘵，何勤如焉？

嘗息肩弛檐於濟岱之閒，或有温故知新雅德君子，矜我劬瘁，睠我皓首，訪論稽古，慰以大道。余困吝之中，精神遐漂，靡所濟集，聊欲係志於翰墨，得以亂思遺老也。

惟六籍之學，先覺之士釋而辯之者，旣已詳矣。儒家惟有《孟子》，閎逺微妙，緼奥難見，宜在條理之科。

於是乃述己所聞，證以經傳，爲之章句，具載本文，章別其指，分爲上下，凡十四卷。究而言之，不敢以當達者。施於新學，可以寤疑辯惑。愚亦未能審於是非，後之明者見其違闕，儻改而正諸，不亦宜乎？


\begin{jieshu}
趙岐《孟子題辭》終
\end{jieshu}

